\documentclass[11pt,a4paper]{article}
\usepackage[T1]{fontenc}
\usepackage[utf8]{inputenc}
\usepackage[english]{babel}
\usepackage{lmodern,amsmath,amssymb,microtype}
\usepackage[margin=25mm]{geometry}
\usepackage[unicode,colorlinks=true,linkcolor=blue,urlcolor=blue]{hyperref}
\hypersetup{pdftitle={Radicals and positive degrees in EGA II},pdfauthor={},pdfsubject={A mathematical note on II, 8.2.10}}
\setlength{\emergencystretch}{2em}
\title{Radicals and positive degrees in EGA II}
\author{}
\date{}
\begin{document}
\maketitle
\subsection*{Radicals and positive degrees in II, 8.2.10}

Let $S=\bigoplus_{d\geq0}S_d$ be a graded ring. Put
$S_{[n]}=\bigoplus_{d\geq n}S_d$ and
$S^\natural=\bigoplus_{n\geq0}S_{[n]}$. In particular,
$(S^\natural)_0=S$. For homogeneous elements $f_\alpha\in S_+$,
write $I=(f_\alpha)\subset S$ and
$J=(f_\alpha^\natural)\subset S^\natural$.
The condition in II, 8.2.10(iv) can be expressed as
\[
 S_+\subseteq\sqrt I,
 \qquad\text{which implies}\qquad
 (S^\natural)_+\subseteq\sqrt J.
\]
Equivalently, intersecting either radical with the corresponding
positive-degree ideal gives that entire ideal. This is the meaning
of the radical within $S_+$ defined in II, 2.1.10. The radical in
the whole ring may contain additional elements of degree zero.

To prove the implication, take $g\in S_{k+h}$ with $k>0$ and
$h\geq0$, viewed as an element of degree $k$ in $S^\natural$.
Such elements additively generate $(S^\natural)_+$.
The hypothesis gives $g^{m_0}=\sum_\alpha c_{\alpha,m_0}f_\alpha$
with only finitely many indices. Multiplication by further powers
of $g$ gives expressions for every $g^m$, $m\geq m_0$, with
the same finite set of indices. Taking homogeneous components,
each nonzero coefficient can be chosen with degree
$m(k+h)-\deg(f_\alpha)$. Choose $m$ sufficiently large that
$km>\deg(f_\alpha)$ for every index in this finite set.
Then $c_{\alpha,m}\in S_{[km-\deg(f_\alpha)]}$ defines a
coefficient of positive degree $km-\deg(f_\alpha)$ in
$S^\natural$. The identity shows that the $m$-th power of $g$,
viewed in degree $k$, belongs to $J$. Since an ideal's radical
is an ideal, the conclusion extends to finite sums.

Equality with the whole-ring radical is false in general.
For a field $k$, let $S=k[e]/(e^2)$, with $\deg(e)=1$ and $f=e$.
Then $S_+=(e)=\sqrt{(e)}$ in $S$, whereas
\[
 S^\natural\simeq k[e_0,e_1]/(e_0^2,e_0e_1,e_1^2),
 \qquad \deg(e_0)=0,\quad\deg(e_1)=1.
\]
Here $f^\natural=e_1$ and $(S^\natural)_+=(e_1)$, but
$\sqrt{(e_1)}=(e_0,e_1)$. The nonzero square-zero element
$e_0$ has degree zero and belongs only to the latter ideal.

\noindent\emph{Reference:} A. Grothendieck, with the collaboration
of J. Dieudonné, \emph{Éléments de géométrie algébrique II},
Publications Mathématiques de l'IHÉS 8 (1961), 2.1.10,
8.2.6 and 8.2.10, pp.~23 and 158--160.


\noindent\href{https://www.numdam.org/item/PMIHÉS_1961__8__5_0/}{French original on NUMDAM}.

\clearpage
\section*{Rights, availability and provenance}
Rights in the original works and in any protected contributions to these derivative translations and transcriptions remain with their respective rights holders. This project does not dedicate those works or contributions to the public domain or grant a licence over them. Existing third-party notices and any permissions already granted remain applicable.

These editions are made available free of charge for non-commercial, purely educational purposes, on the presumption that the respective rights holders would not object to these translations and transcriptions being freely available alongside the French originals on \href{https://www.numdam.org/journals/PMIHÉS/}{NUMDAM}. This presumption is not a statement that permission has been obtained; free access to the originals does not itself grant permission to translate, adapt or republish them.

If you hold rights in any material presented here and object to its availability, please \href{https://github.com/KokunoYumeto/ega-en/issues}{open a GitHub issue}. The publisher will take down the affected material at the rights holder's request.


This mathematical note is AI-created. Its wording, explanation and typesetting were produced by OpenAI Codex, GPT-6 Astra, Ultra. The original mathematical work is by Alexandre Grothendieck, with the collaboration of Jean Dieudonné; it is cited in this note. No approval by the authors, IHÉS or NUMDAM is claimed.
\end{document}
